\newcommand{\DoNotLoadEpstopdf}{}
\pdfinfoomitdate=1
\pdftrailerid{}
\pdfsuppressptexinfo=-1
\documentclass[11pt,letterpaper]{article}
\usepackage[T1]{fontenc}
\usepackage{lmodern}
\usepackage{amsmath,amssymb,amsthm,mathtools}
\usepackage[margin=1in]{geometry}
\usepackage{booktabs,longtable,array}
\usepackage{microtype}
\usepackage[hidelinks]{hyperref}
\hypersetup{pdftitle={Report 157 Airy logarithmic asymptotics for bounded indegree labelled DAGs},pdfauthor={},pdfsubject={OEIS A397711 and every fixed indegree bound}}
\setlength{\parskip}{4pt}
\setlength{\parindent}{1.2em}
\setlength{\emergencystretch}{2em}
\setcounter{tocdepth}{1}
\numberwithin{equation}{section}
\newtheorem{theorem}{Theorem}[section]
\newtheorem{lemma}[theorem]{Lemma}
\newtheorem{proposition}[theorem]{Proposition}
\newtheorem{corollary}[theorem]{Corollary}
\theoremstyle{remark}\newtheorem{remark}[theorem]{Remark}
\newcommand{\NN}{\mathbb N}
\newcommand{\ZZ}{\mathbb Z}
\newcommand{\EE}{\mathbb E}
\newcommand{\PP}{\mathbb P}
\newcommand{\ind}{\mathbf 1}
\newcommand{\floor}[1]{\left\lfloor #1\right\rfloor}
\newcommand{\ceil}[1]{\left\lceil #1\right\rceil}
\DeclareMathOperator{\Ai}{Ai}
\title{Airy logarithmic asymptotics for bounded indegree labelled DAGs\\\large Exact parking reduction and a Lambert W inverse}
\author{Report 157}
\date{3 October 2026}
\begin{document}
\maketitle
\begin{abstract}
Let $a_n$ count simple directed acyclic graphs on the labelled vertex set
$[n]$ with every indegree at most two, the sequence OEIS A397711. If
$-\zeta$ is the largest real zero of the Airy function, then
\[
 \log a_n=2n\log n-n-\kappa n^{1/3}+o(n^{1/3}),
 \qquad \kappa=\frac{3\zeta}{2^{1/3}}.
\]
An exact reduction through generalized parking functions represents the
count as a Poisson excursion partition function with potential
$q_i=\log(i/(i-1))\sim i^{-1}$. Geometric blocks transfer established
homogeneous area-tilt bounds of Mallein to this singular potential.
Explicit interface intervals, a fixed initial segment, and a forced
terminal descent handle the boundaries and the growing number of blocks.

For $y=\log x$ and $b=y/[2W_0(y/(2\sqrt e))]$, the threshold inverse is
\[
 \min\{n\geq1:a_n\geq x\}
 =b+\frac{\kappa b^{1/3}}{2\log b+1}
   +o\!\left(\frac{b^{1/3}}{\log b}\right).
\]
The same argument gives the logarithmic Airy correction for every fixed
integer indegree bound $c\geq2$. Exact enumeration, generalized parking
functions, and the homogeneous Airy estimates are prior inputs. The
result is a logarithmic asymptotic; no power prefactor, amplitude,
$O(\log n)$ remainder, or all-orders expansion is established.
\end{abstract}
\newpage
\tableofcontents
\medskip
\noindent\textbf{Scope and priority.}
All logarithms are natural. The graph model is unweighted and labelled;
graphs are not counted once per compatible topological ordering. The
mathematical statements below follow from the cited probability lemmas,
whose hypotheses are verified explicitly. A bounded literature search
found no equivalent application to this model. This does not establish
first-ever priority. In particular, Steinsky's 2003 paper on bounded
indegree DAG coding was not available for a full-text inspection, so its
priority implications remain unresolved. Section~\ref{sec:sources}
records this limitation separately from the proof.
\newpage

\section{The graph model and main results}\label{sec:main}
A simple directed graph has no loops and at most one edge of each
orientation between two different vertices. A directed acyclic graph
(DAG) has no directed cycle; in particular it cannot contain both
orientations of a pair. Write $a_{n,c}$ for the number of such graphs on
$[n]=\{1,\ldots,n\}$ in which each vertex has at most $c$ incoming edges.
The labels are part of the object. Put $a_{0,c}=1$, and abbreviate
$a_n=a_{n,2}$. The first values are
\[
 a_0,a_1,a_2,a_3,a_4,a_5=1,1,3,25,443,13956.
\]
The sequence is A397711~\cite{oeis}. Exact counting for
fixed indegree bounds is established prior art; see Allen et
al.~\cite{allen}, Theorem~18 and Table~1, and Steinsky~\cite{steinsky}.

Let $\Ai$ be the Airy function and let
\begin{equation}\label{eq:airyconstants}
 \Ai(-\zeta)=0,\qquad
 \zeta=2.338107410459767\ldots,
 \qquad c_A=\frac\zeta{2^{1/3}},\qquad \kappa=3c_A.
\end{equation}
Here $-\zeta$ is the largest real zero, and
$\kappa=5.567271244467717\ldots$. Decimal values are supplied for
orientation; the exact constants are defined by \eqref{eq:airyconstants}.

\begin{theorem}[The first logarithmic correction]\label{thm:main}
As $n\to\infty$,
\begin{equation}\label{eq:main}
 \log a_n=2n\log n-n-\kappa n^{1/3}+o(n^{1/3}).
\end{equation}
Equivalently,
\begin{equation}\label{eq:stretched}
 a_n=\left(\frac{n^2}{e}\right)^n
       \exp\{-\kappa n^{1/3}+o(n^{1/3})\},
 \qquad \lim_{n\to\infty}\frac{a_n^{1/n}}{n^2}=\frac1e.
\end{equation}
\end{theorem}

The error in \eqref{eq:main} is only $o(n^{1/3})$. Formula
\eqref{eq:stretched} is therefore not a ratio equivalent to an explicitly
given sequence. An unproved logarithmic term could still change a power
prefactor, and an unproved constant term could still change an amplitude.
No such refinement is inferred from this theorem or from finite data.

\begin{theorem}[Threshold inversion]\label{thm:inverse}
For real $x\to\infty$, define
\begin{equation}\label{eq:inverseparameters}
 N(x)=\min\{n\geq1:a_n\geq x\},\qquad
 y=\log x,\qquad
 b=\frac{y}{2W_0(y/(2\sqrt e))},
\end{equation}
where $W_0$ is the positive real branch of Lambert's $W$ function on
positive arguments. Then
\begin{equation}\label{eq:inverse}
 N(x)=b+\frac{\kappa b^{1/3}}{2\log b+1}
           +o\!\left(\frac{b^{1/3}}{\log b}\right).
\end{equation}
\end{theorem}

The displacement in \eqref{eq:inverse} is positive. The negative Airy
term lowers the count at a given $n$, so a larger index is needed to
reach $x$. The allowed error scale grows; the theorem supplies
neither eventual exact rounding nor an additive $o(1)$ inverse error.

\begin{theorem}[Every fixed indegree bound]\label{thm:fixedc}
For each fixed integer $c\geq2$, as $n\to\infty$,
\begin{align}
 \log a_{n,c}
 &=cn\log n-\bigl[c-1+\log((c-1)!)\bigr]n\notag\\
 &\hspace{1em}-\kappa(c-1)^{2/3}n^{1/3}+o(n^{1/3}).
 \label{eq:fixedc}
\end{align}
Consequently,
\begin{equation}\label{eq:fixedroot}
 \lim_{n\to\infty}\frac{a_{n,c}^{1/n}}{n^c}
 =\frac{e^{-(c-1)}}{(c-1)!}.
\end{equation}
The error term and its starting index may depend on $c$. No uniformity
is asserted when $c$ grows with $n$.
\end{theorem}

The proof is organized around a potential-transfer lemma for the
Poisson-minus-one walk. The exact counting identities come first so
that the factorials, interval endpoints, and area convention can all be
checked before any limiting argument is applied.

\section{Exact enumeration and generalized parking functions}\label{sec:exact}
\subsection{Inclusion and exclusion over sinks}
Define, with $\binom mj=0$ when $j>m$,
\begin{equation}\label{eq:g}
 g_c(m)=\sum_{j=0}^c\binom mj,\qquad
 g(m)=g_2(m)=1+m+\binom m2=1+\frac{m(m+1)}2.
\end{equation}

\begin{proposition}[Sink recurrence]\label{prop:sink}
For every integer $n\geq1$,
\begin{equation}\label{eq:sink}
 a_{n,c}=\sum_{k=1}^n(-1)^{k+1}\binom nk
                  g_c(n-k)^k a_{n-k,c}.
\end{equation}
\end{proposition}
\begin{proof}
Every nonempty finite DAG has a sink. Fix a set $K$ of $k$ labelled
vertices and require all of them to be sinks. There are no edges inside
$K$ or from $K$ to its complement. The complement supports any allowable
DAG, giving $a_{n-k,c}$ choices. Each sink chooses its parent set
independently from the $n-k$ complementary vertices, with at most $c$
parents, giving $g_c(n-k)^k$ choices. These choices create no directed
cycle. Inclusion and exclusion over the $n$ sink events proves
\eqref{eq:sink}.
\end{proof}

This direct derivation fixes the normalization and model. It does not
assert that bounded-indegree DAG enumeration itself is new. In
particular, one must not multiply the count by a number of topological
orders; that would give a different weighted model.

\subsection{General vector parking bounds}
For a nondecreasing positive integer vector $u_1,\ldots,u_n$, a
$u$-parking preference word is a labelled word of $n$ positive integers
whose sorted values $b_1\leq\cdots\leq b_n$ satisfy $b_i\leq u_i$.
Let $P_n(u_1,\ldots,u_n)$ count these words and put $P_0=1$.
Kung and Yan~\cite{kungyan}, Corollary~5.3 and Theorem~5.4, give the
identity
\begin{equation}\label{eq:parkingidentity}
 x^n=\sum_{j=0}^n\binom nj P_j(u_1,\ldots,u_j)
                         (x-u_{j+1})^{n-j}.
\end{equation}
The last term is $P_n$, so its apparent use of $u_{n+1}$ is harmless.
The counting decomposition initially proves this for every sufficiently
large integer $x$. Both sides are polynomials in $x$, hence it is a
polynomial identity and may be evaluated at $x=0$.

Take
\begin{equation}\label{eq:parkingbounds}
 u_i=g_c(i-1),\qquad u_0=0.
\end{equation}
Moving the $j=n$ term of \eqref{eq:parkingidentity} to the other side and
putting $k=n-j$ gives exactly \eqref{eq:sink}. The initial values agree,
so induction yields
\begin{equation}\label{eq:parkingcount}
 a_{n,c}=P_n\bigl(g_c(0),\ldots,g_c(n-1)\bigr).
\end{equation}
This is an equality of counts obtained from established parking theory;
no new explicit graph-to-parking bijection is required or claimed.

\subsection{Occupancies and the barrier}
Let
\begin{equation}\label{eq:di}
 d_i=u_i-u_{i-1}\quad(i\geq1).
\end{equation}
In the present applications every $d_i$ is positive. For a preference
word, let $x_i$ be the number of its letters in the integer interval
$(u_{i-1},u_i]$. The sorted parking condition is equivalent to
\begin{equation}\label{eq:occupancyconditions}
 x_i\geq0,\qquad \sum_{i=1}^n x_i=n,\qquad
 \sum_{j=1}^i x_j\geq i\quad(1\leq i\leq n).
\end{equation}
Indeed, $b_i\leq u_i$ holds exactly when at least $i$ letters lie at or
below $u_i$. Once the occupancies are fixed, there are
$n!/\prod_i x_i!$ ways to allocate the labelled positions and
$\prod_i d_i^{x_i}$ ways to choose their preference values. Therefore
\begin{equation}\label{eq:occupancy}
 a_{n,c}=n!\sum_{\boldsymbol x\text{ satisfying }\eqref{eq:occupancyconditions}}
                  \prod_{i=1}^n\frac{d_i^{x_i}}{x_i!}.
\end{equation}
For $c=2$,
\begin{equation}\label{eq:quadraticd}
 u_i=1+\binom i2,\qquad d_1=1,\qquad d_i=i-1\quad(i\geq2),
 \qquad \prod_{i=1}^n d_i=(n-1)!.
\end{equation}
The first two increments are both one. This initial duplication will
make $q_1=0$ below; shifting it would change an exact identity.

\section{The exact Poisson excursion representation}\label{sec:poisson}
Let $X_1,X_2,\ldots$ be independent Poisson random variables of mean one,
and define
\begin{equation}\label{eq:walk}
 S_0=0,\qquad S_i=\sum_{j=1}^i(X_j-1),\qquad
 E_n=\{S_i\geq0\ (1\leq i\leq n),\ S_n=0\}.
\end{equation}
Then \eqref{eq:occupancyconditions} is exactly the event $E_n$, under
$x_i=X_i$. For positive $d_i$, put
\begin{equation}\label{eq:qgeneral}
 q_0=0,\qquad q_i=\log(d_{i+1}/d_i)\quad(i\geq1),\qquad
 Z_n(q)=\EE\left[\exp\left\{-\sum_{i=0}^{n-1}q_iS_i\right\};E_n\right].
\end{equation}
The semicolon denotes multiplication by the event indicator.

\begin{proposition}[Exact normalization]\label{prop:normalization}
For every integer $n\geq1$, with the increments \eqref{eq:di} and the
potential \eqref{eq:qgeneral},
\begin{equation}\label{eq:normalization}
 a_{n,c}=n!e^n\left(\prod_{i=1}^n d_i\right)Z_n(q).
\end{equation}
In particular, for $c=2$,
\begin{align}
 q_0=q_1&=0,\qquad q_i=\log\frac{i}{i-1}\quad(i\geq2),\label{eq:q2}\\
 a_n&=n!(n-1)!e^nZ_n(q).\label{eq:normalization2}
\end{align}
\end{proposition}
\begin{proof}
For any fixed occupancy vector, its Poisson probability is
$e^{-n}/\prod_i x_i!$. Equation~\eqref{eq:occupancy} consequently equals
$n!e^n\EE[\prod_i d_i^{X_i};E_n]$. On $E_n$, summation by parts gives
\begin{align}
 \sum_{i=1}^n(X_i-1)\log d_i
 &=\sum_{i=1}^n(S_i-S_{i-1})\log d_i\notag\\
 &=S_n\log d_n-S_0\log d_1
   -\sum_{i=1}^{n-1}S_i\log\frac{d_{i+1}}{d_i}\notag\\
 &=-\sum_{i=1}^{n-1}q_iS_i.\label{eq:summationparts}
\end{align}
Factoring out $\prod_i d_i$ proves \eqref{eq:normalization}.
Equation~\eqref{eq:quadraticd} gives the specialization.
\end{proof}

The area is charged at the left endpoint of each step, from $i=0$ to
$i=n-1$. Its first term vanishes because $S_0=0$. Its last term is
$q_{n-1}S_{n-1}$, not $q_nS_n$. This convention matches the probability
input below and prevents interface costs from being double-counted.
For \eqref{eq:q2}, $q_i\geq0$ and $iq_i\to1$.

\section{The homogeneous Airy estimates used}\label{sec:input}
The external probability input is Mallein~\cite{mallein}, Lemma~2.6,
equation~(2.15), and Lemma~2.7, equation~(2.19), in the arXiv version~3
and corresponding author manuscript. The assumptions are given there
in equations~(2.12)--(2.13). We state precisely the homogeneous
variance-one specialization that is used here.

When the walk starts at $x\in\ZZ_{\geq0}$, write $\EE_x$ for expectation
and still write $S_j$ for its absolute height. For $m\geq1$, $h>0$, and
$J\subseteq[0,\infty)$ define the killed area-tilt kernel
\begin{equation}\label{eq:kernel}
 K_{m,h}(x,J)=\EE_x\left[
 \exp\left\{-\frac hm\sum_{j=0}^{m-1}S_j\right\};
 S_j\geq0\ (0\leq j\leq m),\ S_m\in J\right].
\end{equation}
Every endpoint interval below is intersected with $\ZZ$ when used for
this walk.

\begin{proposition}[Homogeneous bounds from Mallein]\label{prop:input}
For each fixed $h>0$,
\begin{equation}\label{eq:upperinput}
 \limsup_{m\to\infty}m^{-1/3}
  \log\sup_{x\in\ZZ_{\geq0}}K_{m,h}(x,[0,\infty))
 \leq-c_Ah^{2/3}.
\end{equation}
For each fixed $h>0$ and fixed $0<a<b$, $0<a'<b'$,
\begin{equation}\label{eq:lowerinput}
 \liminf_{m\to\infty}m^{-1/3}
 \log\inf_{x\in[a m^{1/3},b m^{1/3}]\cap\ZZ}
 K_{m,h}(x,[a'm^{1/3},b'm^{1/3}])
 \geq-c_Ah^{2/3}.
\end{equation}
\end{proposition}

The source treats independent centered increments with continuous
positive variance profile and a uniform exponential moment of their
absolute values. For $Y=X_1-1$,
\[
 \EE Y=0,\qquad \EE Y^2=1,\qquad
 \EE e^{\mu|Y|}<\infty\quad\text{for every fixed }\mu>0.
\]
Thus the homogeneous array has variance profile identically one and
satisfies every required moment condition. Neither symmetry nor a
nonlattice law is assumed in the cited estimates. The lower estimate
uses an interval of terminal heights, not an exact endpoint. The exact
zero endpoint in our application will be supplied separately.

The normalization of $c_A$ can also be checked at the Brownian level.
The killed operator $\tfrac12 d^2/dx^2-hx$ on the positive half-line
has eigenfunction $\Ai((2h)^{1/3}x-\zeta)$ with eigenvalue
$-\zeta h^{2/3}/2^{1/3}$. This verifies the variance-one convention; it
is not a replacement proof of the random-walk rare-event bounds.

\begin{remark}[The quantifier needed for many blocks]\label{rem:quantifier}
Fix $h$, the endpoint parameters when needed, and $\epsilon>0$.
Equations~\eqref{eq:upperinput}--\eqref{eq:lowerinput} give a single
integer $m_0$ such that the corresponding bounds
\[
 K\leq e^{(-c_Ah^{2/3}+\epsilon)m^{1/3}},\qquad
 K\geq e^{(-c_Ah^{2/3}-\epsilon)m^{1/3}}
\]
hold for every integer $m\geq m_0$, with the indicated supremum or
infimum. This elementary consequence of limsup and liminf is the
uniformity used below. We do not vary $h$ or the scaled endpoint
intervals with the total horizon, and do not require a uniform theorem
over such varying parameters.
\end{remark}

\section{Transferring the estimates to a singular potential}\label{sec:transfer}
\begin{lemma}[A reciprocal potential]\label{lem:transfer}
For the walk \eqref{eq:walk}, let $(q_i)_{i\geq0}$ be finite real numbers
such that $q_i\geq0$ and
\begin{equation}\label{eq:beta}
 iq_i\longrightarrow\beta>0.
\end{equation}
Then
\begin{equation}\label{eq:transfer}
 \log Z_n(q)=-3c_A\beta^{2/3}n^{1/3}+o(n^{1/3}).
\end{equation}
The value of $q_0$ is immaterial because $S_0=0$.
\end{lemma}

The potential is singular after rescaling time near zero. A direct
substitution into a smooth bounded-potential theorem would therefore
need justification. The proof below uses only the homogeneous bounds,
with a fixed initial cutoff and a geometric partition.

\subsection{Geometric grid and rounding control}
Fix $r>1$ and a large constant $T>1$. For integers $N\to\infty$, put
\begin{equation}\label{eq:grid}
 J=\floor{\log_r(N/T)},\qquad
 A_j=Nr^{j-J},\qquad t_j=\floor{A_j}\quad(0\leq j\leq J),
 \qquad m_j=t_j-t_{j-1}.
\end{equation}
In particular, $T\leq A_0<rT$, $T-1\leq t_0<rT$, and $t_J=N$.
Writing $M_j=A_j-A_{j-1}$ gives
\begin{equation}\label{eq:griderrors}
 |m_j-M_j|<1,\qquad M_j=(r-1)A_0r^{j-1}.
\end{equation}
For fixed $r$ and sufficiently large $T$, all $m_j$ are positive and
bounded below by a constant times $Tr^{j-1}$. Uniformly over all the
blocks,
\begin{equation}\label{eq:gridratios}
 \frac{m_j}{t_{j-1}}=(r-1)+O(T^{-1}),\qquad
 \frac{m_j}{t_j}=(1-r^{-1})+O(T^{-1}).
\end{equation}
The constants here may depend on $r$.

The cube-root sum has the limit
\begin{equation}\label{eq:gridsum}
 N^{-1/3}\sum_{j=1}^Jm_j^{1/3}\longrightarrow
 D(r)=\frac{(1-r^{-1})^{1/3}}{1-r^{-1/3}}.
\end{equation}
To verify rounding does not affect this limit, the mean value theorem
and \eqref{eq:griderrors} give
\[
 |m_j^{1/3}-M_j^{1/3}|\leq C_rM_j^{-2/3}
\]
when $T$ is sufficiently large. The right side is summable over the
geometric grid, with a sum bounded independently of $N$. The exact
geometric sum is
\[
 N^{-1/3}\sum_{j=1}^JM_j^{1/3}
 =(1-r^{-1})^{1/3}\sum_{k=0}^{J-1}r^{-k/3},
\]
which tends to $D(r)$. This proves \eqref{eq:gridsum}.

The minimum block length is not asserted to tend to infinity with $N$.
Instead, $T$ is chosen sufficiently large once the fixed parameters and
required length thresholds have been selected. It then remains fixed
while $N\to\infty$.

\subsection{Comparing the potential on each block}
Fix also $\rho\in(0,1)$ and set
\begin{equation}\label{eq:hpm}
 h_-=(1-\rho)\beta(1-r^{-1}),\qquad
 h_+=(1+\rho)\beta(r-1).
\end{equation}
After increasing $T$, condition~\eqref{eq:beta} and
\eqref{eq:gridratios} imply simultaneously on every block
\begin{equation}\label{eq:comparison}
 \frac{h_-}{m_j}\leq q_i\leq\frac{h_+}{m_j}
 \qquad(t_{j-1}\leq i<t_j).
\end{equation}
For example, first bound $iq_i$ between
$(1-\delta)\beta$ and $(1+\delta)\beta$ for all indices above a fixed
cutoff. On the block use $i\leq t_j$ for the lower bound and
$i\geq t_{j-1}$ for the upper bound. Choose $\delta>0$ small in terms
of $r,\rho$, then enlarge $T$ to absorb the ratio errors. No monotonicity
of $q_i$ is needed. Nonnegativity permits the kernel comparisons on
nonnegative paths and removal of the initial cost in the upper bound.

\subsection{Upper bound}
Take $N=n$. Remove the final requirement $S_n=0$ and the nonnegative
area cost before $t_0$. The remaining initial distribution on nonnegative
heights has total mass at most one. The Markov property and
\eqref{eq:comparison} yield
\begin{equation}\label{eq:upperproduct}
 Z_n(q)\leq\prod_{j=1}^J
      \sup_{x\in\ZZ_{\geq0}}K_{m_j,h_-}(x,[0,\infty)).
\end{equation}
For fixed $r,\rho,\epsilon>0$, choose $T$ still larger so that every
$m_j$ exceeds the single threshold in
Remark~\ref{rem:quantifier}. Consequently,
\[
 \log Z_n(q)\leq(-c_Ah_-^{2/3}+\epsilon)
                       \sum_{j=1}^Jm_j^{1/3}.
\]
Take $n\to\infty$, then $\epsilon\downarrow0$, then
$\rho\downarrow0$, while $r$ remains fixed. Formula~\eqref{eq:gridsum}
and the algebraic identity
\begin{equation}\label{eq:upperconstant}
 (1-r^{-1})^{2/3}D(r)=1+r^{-1/3}+r^{-2/3}
\end{equation}
give
\begin{equation}\label{eq:upperfinal}
 \limsup_{n\to\infty}n^{-1/3}\log Z_n(q)
 \leq-c_A\beta^{2/3}(1+r^{-1/3}+r^{-2/3}).
\end{equation}

\subsection{Lower bound interfaces with explicit slack}
Fix in addition $\eta>0$, and reserve a short terminal interval by
setting
\begin{equation}\label{eq:tail}
 L_n=\ceil{2\eta n^{1/3}},\qquad N=n-L_n.
\end{equation}
Use \eqref{eq:grid} with this $N$. At ordinary interfaces require
\begin{equation}\label{eq:interfaces}
 S_{t_j}\in I_j=[t_j^{1/3},2t_j^{1/3}]
 \qquad(0\leq j<J),
\end{equation}
and at the last interface require
\begin{equation}\label{eq:lastinterface}
 S_N\in I_J^{(\eta)}=[\eta N^{1/3},2\eta N^{1/3}].
\end{equation}
Only sufficiently large $n$, hence $J\geq1$, is relevant.

The following fixed intervals make the choice of parameters transparent.
Put
\begin{equation}\label{eq:alphagamma}
 \alpha=(r-1)^{-1/3},\qquad
 \gamma=\left(\frac r{r-1}\right)^{1/3}.
\end{equation}
By \eqref{eq:gridratios}, after increasing $T$ the scaled incoming
sets $I_{j-1}/m_j^{1/3}$ all lie in
\begin{equation}\label{eq:incominginterval}
 [a,b]=[\alpha/2,3\alpha].
\end{equation}
All ordinary scaled outgoing sets $I_j/m_j^{1/3}$ contain the fixed
interval
\begin{equation}\label{eq:outgoinginterval}
 [a',b']=[5\gamma/4,7\gamma/4],
\end{equation}
and the last set $I_J^{(\eta)}/m_J^{1/3}$ contains
\begin{equation}\label{eq:lastoutgoing}
 [a'_\eta,b'_\eta]=[5\eta\gamma/4,7\eta\gamma/4].
\end{equation}
Indeed, the limiting ordinary outgoing interval is $[\gamma,2\gamma]$,
and the last one is $[\eta\gamma,2\eta\gamma]$, so these choices have
strict positive slack. These intervals depend on the fixed $r,\eta$
only. We can therefore select the two lower-bound length thresholds
first and subsequently increase $T$ to exceed their maximum. Increasing
$T$ preserves the required containments, so there is no circular choice.

\subsection{A finite initial segment}
For any possible $t_0$, choose an integer
$H\in[t_0^{1/3},2t_0^{1/3}]$. Force $X_1=H+1$ and
$X_2=\cdots=X_{t_0}=1$. This path is nonnegative, ends at height $H$,
and has weighted probability
\begin{equation}\label{eq:initialcost}
 \frac{e^{-t_0}}{(H+1)!}
       \exp\left\{-H\sum_{i=1}^{t_0-1}q_i\right\}.
\end{equation}
For fixed $T,r$, both $t_0$ and the selected $H$ take values in finite
sets. Because each $q_i$ is finite, the minimum of
\eqref{eq:initialcost} over these possibilities is a positive constant
$C_{T,r,q}$. It may be very small as $T$ increases, but $T$ remains fixed
in the limit in $n$, so $n^{-1/3}\log C_{T,r,q}\to0$.

\subsection{Multiplying the lower kernels}
Apply \eqref{eq:lowerinput} at $h=h_+$, with incoming interval
\eqref{eq:incominginterval}, and with the ordinary or last outgoing
interval as appropriate. For fixed $r,\rho,\eta,\epsilon$, enlarge
$T$ so all lengths exceed both thresholds. The actual block transition
mass is at least
\begin{equation}\label{eq:lowerfactor}
 \exp\{(-c_Ah_+^{2/3}-\epsilon)m_j^{1/3}\}.
\end{equation}
The reason is that \eqref{eq:comparison} bounds the actual potential
above by $h_+/m_j$, so its weight is at least the homogeneous weight.

For a given incoming height in $I_{j-1}$, the selected outgoing interval
is contained in $I_j$ (or in $I_J^{(\eta)}$ for the last block). Every
resulting point of $I_j$ is covered by the incoming infimum for the next
block. The Markov property thus multiplies \eqref{eq:lowerfactor} without
an endpoint matching loss. Each transition mass already sums over its
terminal lattice points; there is no division by an interval length or
additional per-state factor. The left-endpoint area convention assigns
each interface cost exactly once, to its outgoing block.

Although the number of blocks grows with $n$, the total logarithmic
error in their product is only
$\epsilon\sum_jm_j^{1/3}$. After division by $N^{1/3}$ it tends to
$\epsilon D(r)$. This uses the all-$m\geq m_0$ bound, not an unjustified
exchange of a blockwise limit with an increasing product.

\subsection{Closing the excursion at zero}
At time $N$ let the integer height be
$s\in I_J^{(\eta)}$. Then
$s\leq2\eta N^{1/3}\leq L_n$. For the remaining $L_n$ steps force
$X_i=0$ for the first $s$ steps and $X_i=1$ for the remaining steps.
The walk descends by one to zero, then stays at zero. Both prescribed
outcomes have probability $e^{-1}$, so this conditional event has
probability exactly $e^{-L_n}$, independently of $s$.

Throughout the tail the height is at most $L_n$. Since $N/n\to1$,
condition~\eqref{eq:beta} gives
$\max_{N\leq i<n}q_i=O(n^{-1})$. Therefore its extra area satisfies
\begin{equation}\label{eq:tailarea}
 \sum_{i=N}^{n-1}q_iS_i
 \leq L_n^2\max_{N\leq i<n}q_i=O(n^{-1/3}).
\end{equation}
The weighted tail probability is at least
$\exp\{-L_n-O(n^{-1/3})\}$. This step uses the special support of the
Poisson-minus-one walk; Lemma~\ref{lem:transfer} is asserted precisely
for that walk, which is also the walk used for all fixed $c$.

Combining the initial segment, all blocks, and the terminal tail gives
\begin{align}
 \log Z_n(q)\geq\log C_{T,r,q}
   +(-c_Ah_+^{2/3}-\epsilon)\sum_{j=1}^Jm_j^{1/3}
   -L_n-O(n^{-1/3}).\label{eq:lowerproduct}
\end{align}
First take $n\to\infty$. Then send $\epsilon,\rho\downarrow0$, with
$r,\eta$ fixed. Using $N/n\to1$, \eqref{eq:gridsum}, and
\begin{equation}\label{eq:lowerconstant}
 (r-1)^{2/3}D(r)=1+r^{1/3}+r^{2/3},
\end{equation}
we obtain
\begin{equation}\label{eq:lowerfinal}
 \liminf_{n\to\infty}n^{-1/3}\log Z_n(q)
 \geq-c_A\beta^{2/3}(1+r^{1/3}+r^{2/3})-2\eta.
\end{equation}
Finally send $\eta\downarrow0$ and $r\downarrow1$ in
\eqref{eq:upperfinal} and \eqref{eq:lowerfinal}. Both bounds tend to
$-3c_A\beta^{2/3}$, proving Lemma~\ref{lem:transfer}.

\begin{remark}[Order of limits]
The fixed parameters $r,\rho,\eta,\epsilon$ are selected first, then
$T$ is fixed sufficiently large, and only then $n$ tends to infinity.
At each fixed $r$, the error parameter $\epsilon$ tends to zero before
$r\downarrow1$. This matters because $D(r)$ diverges as $r\downarrow1$.
No quantitative rate in the homogeneous estimates is needed or claimed.
\end{remark}

\section{Completion of the DAG asymptotics}\label{sec:completion}
For $c=2$, the exact potential \eqref{eq:q2} satisfies
$iq_i\to1$. Lemma~\ref{lem:transfer} therefore gives
\begin{equation}\label{eq:zlimit}
 \log Z_n(q)=-\kappa n^{1/3}+o(n^{1/3}).
\end{equation}
Using $(n-1)!=n!/n$, Stirling's formula yields
\begin{align}
 \log\bigl(n!(n-1)!e^n\bigr)
 &=2\log(n!)-\log n+n\notag\\
 &=2n\log n-n+\log(2\pi)+O(n^{-1}).\label{eq:stirling}
\end{align}
Combining \eqref{eq:normalization2}, \eqref{eq:zlimit}, and
\eqref{eq:stirling} proves Theorem~\ref{thm:main}.

The two factorials cancel their logarithmic terms. This cancellation
is an exact normalization check, not evidence of a missing universal
prefactor. Any logarithmic contribution in $\log a_n$ beyond
\eqref{eq:stirling} must come from a sharper analysis of $Z_n(q)$;
Lemma~\ref{lem:transfer} does not provide that sharper analysis.

\subsection{Extension to every fixed c}
Fix an integer $c\geq2$. Pascal's identity applied to \eqref{eq:g} gives
\begin{equation}\label{eq:generald}
 d_1=1,\qquad
 d_i=g_c(i-1)-g_c(i-2)
     =\sum_{j=0}^{c-1}\binom{i-2}{j}\quad(i\geq2).
\end{equation}
These increments are positive and nondecreasing, so
$q_i=\log(d_{i+1}/d_i)\geq0$. As $i\to\infty$ with $c$ fixed,
\begin{equation}\label{eq:dasymptotic}
 d_i=\frac{i^{c-1}}{(c-1)!}\bigl(1+O(i^{-1})\bigr),
 \qquad iq_i\longrightarrow c-1.
\end{equation}
For the second assertion one may expand the fixed polynomial one term
further: $d_i=A i^{c-1}(1+B/i+O(i^{-2}))$ for constants $A>0,B$.
Taking a ratio at consecutive integers gives
$q_i=(c-1)/i+O(i^{-2})$.

Summing logarithms in \eqref{eq:dasymptotic} gives
\begin{equation}\label{eq:productgeneral}
 \log\prod_{i=1}^n d_i
 =(c-1)\log(n!)-n\log((c-1)!)+O(\log n).
\end{equation}
Finite initial indices are absorbed in the error. Proposition
\ref{prop:normalization}, Lemma~\ref{lem:transfer} with $\beta=c-1$,
and Stirling's formula now give
\begin{align*}
 \log a_{n,c}
 &=c\log(n!)+n-n\log((c-1)!)
   -\kappa(c-1)^{2/3}n^{1/3}+o(n^{1/3})+O(\log n)\\
 &=cn\log n-\bigl[c-1+\log((c-1)!)\bigr]n
   -\kappa(c-1)^{2/3}n^{1/3}+o(n^{1/3}).
\end{align*}
This proves Theorem~\ref{thm:fixedc}. The $O(\log n)$ error in the
separate product calculation is absorbed by $o(n^{1/3})$; it does not
improve the final remainder to $O(\log n)$.

For context, $c=1$ is a different boundary case: such DAGs are rooted
labelled forests with edges oriented away from the roots, and
$a_{n,1}=(n+1)^{n-1}$ for $n\geq1$. In that case $d_i=1$ and $q_i=0$,
so $\beta=0$. The positive-$\beta$ transfer lemma is not used to claim
an Airy correction there. The forest formula is a useful independent
finite check of the exact recurrence.

\section{Lambert W and the threshold inverse}\label{sec:inverse}
The threshold in \eqref{eq:inverseparameters} exists because the counts
grow without bound. They are strictly increasing for $n\geq2$: extend
any allowed graph on $[n-1]$ by vertex $n$ as a sink, either isolated or
with one of the $n-1$ old vertices as its sole parent. These $n$
extensions are distinct and respect the indegree bound, so
\begin{equation}\label{eq:monotonicity}
 a_n\geq n a_{n-1}\qquad(n\geq2).
\end{equation}

Put $F(t)=2t\log t-t$. If
$w=W_0(y/(2\sqrt e))$, then $we^w=y/(2\sqrt e)$ and
\[
 b=\frac y{2w}=\sqrt e\,e^w,\qquad
 \log b=\frac12+w,\qquad F(b)=2bw=y.
\]
Thus $b$ is the exact large positive solution of the leading smooth
model, and
\begin{equation}\label{eq:derivatives}
 F'(b)=2\log b+1,\qquad F''(b)=2/b.
\end{equation}

Here are explicit bracketing details for the error precision.
Fix $\epsilon>0$. Theorem~\ref{thm:main} implies that for every
sufficiently large integer $n$,
\begin{equation}\label{eq:countbrackets}
 F(n)-(\kappa+\epsilon)n^{1/3}
 \leq\log a_n\leq
 F(n)-(\kappa-\epsilon)n^{1/3}.
\end{equation}
For any fixed real $A$, let
\[
 t_A=b+\frac{A b^{1/3}}{F'(b)}.
\]
Taylor's theorem and \eqref{eq:derivatives} give
\begin{equation}\label{eq:taylorinverse}
 F(t_A)=y+A b^{1/3}
        +O\!\left(\frac{b^{-1/3}}{\log^2 b}\right),
 \qquad t_A^{1/3}=b^{1/3}
        +O\!\left(\frac{b^{-1/3}}{\log b}\right).
\end{equation}
Integer rounding changes $F(t_A)$ by $O(\log b)$, which is
$o(b^{1/3})$. Choosing $A=\kappa-2\epsilon$ and
$A=\kappa+2\epsilon$ in \eqref{eq:countbrackets} shows, for all
sufficiently large $x$,
\begin{equation}\label{eq:inversebrackets}
 \floor{t_{\kappa-2\epsilon}}<N(x)
 \leq\ceil{t_{\kappa+2\epsilon}}.
\end{equation}
For the left inequality the upper count bound is below $y$; for the
right inequality the lower count bound is above $y$. The strict
margins are of order $\epsilon b^{1/3}$ and dominate the Taylor and
rounding errors.

Finally, $O(1)=o(b^{1/3}/\log b)$. Divide the displacement bounds by
$b^{1/3}/F'(b)$ and let $\epsilon\downarrow0$. This proves
\eqref{eq:inverse}. No numerical asymptotic onset or explicit
remainder constant is supplied by this proof, so it does not certify a
finite-input numerical bracket for $N(x)$.

\section{Exact companion and reproducibility}\label{sec:companion}
The accompanying source package contains the editable LaTeX article,
the PDF, a Python standard-library companion, tests, attributed numerical
prefixes, deterministic exact output, source provenance, and controlled
build and archive helpers. No journal full text or unrelated research
material is bundled.

The count calculation uses the sink recurrence
\eqref{eq:sink} with arbitrary-precision integers. Independently
structured checks use the parking occupancy constraints, rational
Poisson area weights, and direct small-graph enumeration. In the latter,
each unordered vertex pair has three possible states: absent, one
orientation, or the other. An indegree check and an acyclicity test then
select precisely the required labelled simple DAGs. Small parking-word
enumeration checks the weak upper-bound convention directly. The
$c=1$ forest formula supplies another independent reference case.

The exact threshold operation searches a bounded table of actual
counts. If its chosen range is too short, it reports that fact rather
than replacing an exact answer by the asymptotic inverse. Workload
limits are software limits only and do not restrict the mathematical
theorems. Finite checks support the formulas' indexing and the code's
correctness; they prove neither an asymptotic rate nor scholarly
priority.

The README records the exact command lines, bounds, shipped test
coverage, and output policies. Validation is explicit rather than
implemented with Python assertions that disappear under optimization;
the suite is run both normally and with \texttt{-O}. Output is
exclusive and refuses existing targets and symbolic-link traversal.
These are local safety measures under a trusted-parent and
nonadversarial-concurrent-environment assumption, not an operating-system
sandbox or a claim of general race resistance.

For the PDF, the helper creates an exclusive clean build directory,
uses a fresh LaTeX format and cache, disables shell escape, fixes PDF
metadata, and repeats compilation until references stabilize. The
settled log must contain no warnings or layout diagnostics. Identical
bytes are expected on repeated builds within the recorded toolchain,
not across arbitrary TeX versions, fonts, or platforms. In particular,
the directory creation and subsequent opening require a trusted parent
and no adversarial concurrent replacement.

The release manifest covers a fixed sorted payload allowlist. The ZIP
uses fixed timestamps, permissions, order, and stored compression; it
includes the checksum manifest, which necessarily excludes its own
self-hash. A separately reported archive digest identifies the complete
release. The supplied commands permit checksum validation, fresh
extraction, exact-output replay, normal and optimized tests, clean PDF
rebuild, and deterministic archive replay.

\section{Attribution and limits of source clearance}\label{sec:sources}
The argument has three distinct prior inputs. First, bounded-indegree
labelled DAG enumeration is already treated by Steinsky~\cite{steinsky}
and Allen et al.~\cite{allen}. Second, the general-vector parking
identity is due to Kung and Yan~\cite{kungyan}; Yan's survey~\cite{yan}
provides additional context. Third, the homogeneous killed-walk
area-tilt estimates are Mallein's~\cite{mallein}. The proof here applies
these inputs to a reciprocal potential, verifies the endpoint and
block quantifiers, and derives the stated DAG logarithmic asymptotic
and threshold displacement. It does not claim any of the underlying
mechanisms as new.

The source preparation for this report inspected the relevant exact
counting section of Allen et al., all pages of the Kung--Yan author
manuscript through rendering and OCR, and the specific Mallein
assumptions, lemmas, and nearby proofs used above. The crucial formula
pages were checked visually because the parking manuscript's original
text extraction was damaged. The homogeneous estimates also agree
between Mallein's arXiv version~3 and author manuscript. These statements
record the scope of the source check, not a claim that all surrounding
literature was exhaustively reviewed.

Steinsky's publisher abstract verifies bounded-indegree enumeration,
generation, ranking, and unranking, but the full text was not read in
this source pass. No assertion is made that it contains no relevant
asymptotic. This unresolved full-text priority check remains important
before any claim of novelty or publication priority.

Related models require care. Chen and Pearl's bounded-indegree
Bayesian-network sampling model~\cite{chenpearl} weights DAGs through
compatible orders, whereas the present count gives every labelled graph
weight one. General parking-function asymptotics with affine bounds,
such as Yin~\cite{yin}, concern a different boundary from the quadratic
$u_i=1+\binom i2$ here. Results for related graph and parking models do not by themselves
identify the enumeration of arbitrary labelled simple DAGs. A bounded search found no equivalent
result for this exact model, but absence from a search is not a proof
of global priority.

The remaining analytic target is a quantitative expansion of the
inhomogeneous excursion partition function that controls terms below
$n^{1/3}$. Such a result would be needed to isolate a logarithmic
prefactor, an amplitude, or further powers. None is supplied here.
Similarly, no uniform extension to growing $c$ is obtained: all constants,
cutoffs, and probability parameters in the fixed-$c$ argument are chosen
after $c$ is fixed.

\begin{thebibliography}{99}
\raggedright
\bibitem{allen}
T.~E. Allen, J. Goldsmith, H.~E. Justice, N. Mattei, and K. Raines,
\emph{Uniform Random Generation and Dominance Testing for CP-Nets},
Journal of Artificial Intelligence Research \textbf{59} (2017), 771--813.
Theorem~18 and Table~1.
\url{https://par.nsf.gov/servlets/purl/10048706}.

\bibitem{steinsky}
B. Steinsky, \emph{Efficient coding of labeled directed acyclic graphs},
Soft Computing \textbf{7} (2003), 350--356.
\href{https://doi.org/10.1007/s00500-002-0223-5}{doi:10.1007/s00500-002-0223-5}.
Publisher abstract inspected; full text not inspected.

\bibitem{oeis}
OEIS Foundation Inc., \emph{A397711}, number of labelled simple DAGs
with each indegree at most two. Official data mirror consulted
3 October 2026; entry version header 12 July 2026.
\url{https://oeis.org/A397711}.
\url{https://github.com/oeis/oeisdata/blob/main/seq/A397/A397711.seq}.

\bibitem{kungyan}
J.~P.~S. Kung and C.~H. Yan, \emph{Gon\v{c}arov polynomials and parking
functions}, Journal of Combinatorial Theory, Series A \textbf{102}
(2003), 16--37. Corollary~5.3 and Theorem~5.4.
\href{https://doi.org/10.1016/S0097-3165(03)00009-8}{doi:10.1016/S0097-3165(03)00009-8}.
Author manuscript: \url{https://people.tamu.edu/~huafei-yan/Files/GP.pdf}.

\bibitem{yan}
C.~H. Yan, \emph{Parking Functions}, in M. B\'ona (ed.),
Handbook of Enumerative Combinatorics, CRC Press (2015), 835--893.
Author copy:
\url{https://people.tamu.edu/~huafei-yan/Files/Yan-Final-Own-Copy.pdf}.

\bibitem{mallein}
B. Mallein, \emph{Maximal displacement of a branching random walk in
time-inhomogeneous environment}, Stochastic Processes and their
Applications \textbf{125} (2015), 3958--4019.
\href{https://doi.org/10.1016/j.spa.2015.05.011}{doi:10.1016/j.spa.2015.05.011}.
\url{https://arxiv.org/abs/1307.4496}; version~3, 19 May 2015,
assumptions~(2.12)--(2.13), Lemmas~2.6--2.7, equations~(2.15) and~(2.19).
Author manuscript:
\url{https://www.math.univ-toulouse.fr/~bmallein/publications/brwtie.pdf}.

\bibitem{chenpearl}
E.~Y.-J. Chen and J. Pearl, \emph{Random Bayesian networks with bounded
indegree}, Proceedings of AISTATS 2014, Proceedings of Machine Learning
Research \textbf{33} (2014).
\url{https://proceedings.mlr.press/v33/chen14a.pdf}.

\bibitem{yin}
M. Yin, \emph{Parking Functions, Multi-Shuffle, and Asymptotic Phenomena},
33rd International Conference on Probabilistic, Combinatorial and
Asymptotic Methods for the Analysis of Algorithms (AofA 2022),
Leibniz International Proceedings in Informatics \textbf{225},
Article~18.
\href{https://doi.org/10.4230/LIPIcs.AofA.2022.18}{doi:10.4230/LIPIcs.AofA.2022.18}.
\end{thebibliography}
\end{document}
